\section{What the design could have found}\label{sec:power}

A null result is informative only in proportion to the smallest effect it could
have detected. For a partial effect estimated on $n$ observations,
$t^2 \approx n\,R^2/(1-R^2)$, so at a fixed bar $t^\ast$ the smallest detectable
\pRsq{} is
\begin{equation}
  R^2_{\min} = \frac{t^{\ast 2}}{n + t^{\ast 2}},
  \label{eq:power}
\end{equation}
which falls as $1/n$. With \powerN{} discovery weeks and the Bonferroni bar of
$|t|>\bonferroniT$, the floor is \powerFloor{} of residual variance
(Figure~\ref{fig:power}). The largest \pRsq{} observed anywhere in the paper, across
eight signals, two horizons, three mechanisms and two samples, is \bMaxRsqCov.

\begin{figure}[t]
\centering
\includegraphics[width=\linewidth]{fig4_power}
\caption{Smallest detectable \pRsq{}, Equation~\eqref{eq:power}, as a function of
sample size at the two bars used in the paper, with the largest observed \pRsq{} and
the discovery sample size marked. Halving the floor needs twice the sample, which
at a weekly grid means another six years of coverage-consistent data.}
\label{fig:power}
\end{figure}

The scope of the result follows from this. It excludes effects large enough to be
traded and is silent on smaller ones. A signal explaining a third of a percent of
residual variance would be physically real and commercially useless, and this design
could not distinguish it from zero. Two facts about scale make the null less
surprising than the intuition that motivated the study. The EU-bound stock of gas at
sea averages \stockMeanMm{} million\,m$^3$ of LNG, which is about
\stockMeanBcm{}\,bcm of gas or \stockDaysDemand{} days of EU demand at roughly
\euDemandBcm{}\,bcm per year. A one-standard-deviation change in that stock is
\stockSdDaysDemand{} days of demand. The weekly spread has a standard deviation of
\$\spreadWeeklySd{} per MMBtu. The question was therefore whether about a day and a
half of marginal supply moves a price that responds to revisions in the weather
forecast.

\paragraph{Distinguishing a physically real signal from a tradeable one.} The
distinction is worth keeping because it is where the interesting work now lies. A
signal is physically real if it predicts the quantity it measures: Part~A is that
test, and the at-sea stock does predict European arrivals, with a lead equal to the
voyage time, because that is arithmetic. A signal is tradeable if its effect on the
price survives the controls at a size above the floor of
Equation~\eqref{eq:power}. The route to the first question is more physical
targets scored against published totals; the route to the second is a lower floor,
which means a longer coverage-consistent sample or a less attenuated regressor, not
a different model.
